% 第四章 重点板块深度分析（一）：高确定性预期差

\chapter{重点板块深度分析（一）：高确定性预期差}

本章深入分析两大高确定性预期差板块——功率半导体与电力公用事业。两大板块均已通过"业绩验证—产业映射—政策加持—估值持仓"四维度交叉验证，景气度真实且市场预期尚未充分定价，是当前市场环境下胜率与赔率兼备的配置方向。

\section{功率半导体：AI算力+能源革命双轮驱动，国产替代纵深推进}

功率半导体是电力电子系统的"心脏"，广泛应用于新能源车、工业控制、光伏储能、数据中心等领域。我们认为，功率半导体正处于新一轮上行周期的起步阶段，AI算力革命与能源转型两大超级周期共振，叠加国产替代从消费级向车规级纵深推进，行业成长性显著强于周期性，是半导体板块中最具配置价值的细分赛道之一。

\subsection{景气度三维验证}

\textbf{业绩维度：行业底部确认，复苏信号明确}。经过2023--2024年的深度调整，功率半导体行业库存去化基本完成，主要厂商存货周转天数从峰值的180天回落至120天左右的正常水平。2025年Q1以来，行业营收同比增速逐季改善，主要龙头公司毛利率企稳回升。根据Wind一致预期，2025年A股功率半导体板块归母净利润同比增速约为25--30\%，2026年有望达到35--40\%，业绩拐点已基本确立。

\vspace{0.3em}

\textbf{产业维度：AI+新能源双轮驱动，需求多点开花}。AI服务器是功率半导体边际增量最大的赛道。从传统CPU服务器到AI GPU服务器，单机架功耗从300--500W飙升至10--20kW，功率半导体价值量提升8--12倍。下一代Blackwell架构服务器单机功率半导体价值量可达5--6万美元，较上一代提升约50\%。我们测算，全球AI服务器功率半导体市场规模将从2024年的约39亿美元增长至2029年的约98亿美元，CAGR约20\%，显著高于整体功率半导体市场5--6\%的增速。

\begin{exhibitbox}[表4-1：服务器演进与功率半导体价值量对比]
\centering
\small
\begin{tabular}{l r r r}
\toprule
\textbf{服务器类型} & \textbf{单机架功率} & \textbf{功率半导体价值量} & \textbf{相对倍数} \\
\midrule
传统机架式服务器 & 0.3--0.5 kW & \$2,000--5,000 & 1$\times$ \\
AI训练服务器（8-GPU） & 10--12 kW & \$30,000--40,000 & 8--12$\times$ \\
下一代AI服务器（B100） & 15--20 kW & \$50,000--60,000 & 12--20$\times$ \\
超高密度整机柜（GB200） & 70--80 kW & \$70,000--80,000 & 15--25$\times$ \\
\bottomrule
\end{tabular}
\sourcenote{Nvidia，本研究团队测算（按整机柜功率半导体价值估算）}
\end{exhibitbox}

新能源车是功率半导体第一大下游，约占35--40\%的需求。2025年国内新能源车销量预计达1,300--1,400万辆，同比增长约25\%，单车功率半导体价值量随800V平台和SiC渗透持续提升。此外，光伏储能、工业控制、家电变频等传统需求也呈现稳健增长态势。

\vspace{0.3em}

\textbf{政策维度：国产替代战略地位提升，政策支持加码}。半导体产业政策持续加码，大基金三期、税收优惠、车规芯片扶持等政策密集出台。功率半导体作为半导体领域国产替代最确定、进度最快的赛道之一，政策红利持续释放。特别是在车规级芯片领域，自主可控的战略意义日益凸显，政策支持力度有望进一步加大。

\subsection{预期差来源分析}

我们认为，当前市场对功率半导体的认知存在三方面显著预期差：

\vspace{0.3em}

\textbf{预期差一：低估AI侧增量弹性}。市场普遍将功率半导体视为"新能源+周期"板块，对AI带来的增量弹性认识不足。多数投资者认为AI对功率半导体的拉动有限且集中在海外厂商，但实际上，AI服务器电源架构从12V向48V演进、VRM相数翻倍、GaN器件加速渗透，均为国内厂商带来新的增长机遇。随着国内电源芯片厂商逐步进入AI服务器供应链，国产替代空间巨大。我们测算，AI侧需求对国内功率半导体厂商的收入贡献将从2024年的不足5\%提升至2027年的15--20\%，成为第二增长曲线。

\vspace{0.3em}

\textbf{预期差二：对SiC/GaN渗透率预期偏保守}。市场对第三代半导体的渗透率预期普遍偏保守，多数预期认为2027年SiC在新能源车中的渗透率约20--25\%。但我们认为，随着800V平台加速上车、SiC成本快速下降（预计2026--2027年SiC器件价格降至硅基IGBT的1.5--2倍），渗透率有望超预期。我们判断，2027年国内新能源车SiC渗透率有望达到30--35\%，2030年有望突破50\%。GaN在数据中心VRM和快充领域的渗透速度同样可能超预期。

\vspace{0.3em}

\textbf{预期差三：国产替代深度与速度被低估}。市场普遍认为功率半导体国产替代主要集中在中低端消费级和工业级领域，车规级仍需较长时间。但实际上，国内头部厂商车规级IGBT已进入大规模量产阶段，部分车企国产供应商份额已达20--30\%，且提升速度正在加快。SiC领域，国内厂商虽起步较晚，但衬底良率提升和器件验证进度快于市场预期。我们判断，2027年国内车规级功率半导体国产化率有望从当前的约15\%提升至30\%以上。

\subsection{核心催化与时间窗口}

未来12--18个月，功率半导体板块存在五大核心催化，按时间先后排序：

\vspace{0.3em}

\textbf{催化一：Blackwell架构渗透加速（2026H2--2027H1）}。英伟达Blackwell架构B100/B200芯片大规模出货，单机功率和功率半导体价值量进一步提升，带动上游功率器件需求。关注指标：英伟达出货量、AI服务器出货量、电源厂商订单情况。

\vspace{0.3em}

\textbf{催化二：800V平台上车加速（2026H2）}。比亚迪、吉利、小鹏、理想等车企800V平台车型密集推出，带动SiC器件需求放量。关注指标：800V车型销量、SiC器件出货量、车企平台化战略进展。

\vspace{0.3em}

\textbf{催化三：国产替代从消费级迈向车规级（2026Q3--2027Q1）}。国内厂商车规级IGBT和SiC器件从验证走向批量装车，份额快速提升。关注指标：车规业务收入占比、AEC-Q101认证进展、头部车企供应链份额。

\vspace{0.3em}

\textbf{催化四：行业价格企稳回升（2026Q3--Q4）}。随着需求复苏和库存去化完成，功率半导体产品价格有望企稳回升，带来业绩弹性。关注指标：渠道价格走势、晶圆代工产能利用率、厂商毛利率变化。

\vspace{0.3em}

\textbf{催化五：大基金三期投资落地（2026下半年）}。大基金三期对半导体设备材料和功率器件的投资落地，提振市场情绪。关注指标：大基金投资进展、产业政策出台。

\subsection{产业链图谱与受益环节}

功率半导体产业链可分为上游材料设备、中游制造、下游应用三个环节。从投资弹性看，中游制造环节受益最直接，上游材料设备次之。

\begin{exhibitbox}[图4-1：功率半导体产业链与受益环节图谱]
\centering
\vspace{0.4cm}
\makebox[\textwidth][c]{%
\begin{tikzpicture}[node distance=0.4cm, every node/.style={font=\small}]

% ===== 上游：材料与设备 =====
\node[draw, fill=navy!20, rectangle, rounded corners=2pt, minimum height=0.85cm, minimum width=2.0cm, align=center] (wafer) at (-4.8,2.0) {\textbf{硅片/衬底}\\ 沪硅产业、立昂微};
\node[draw, fill=navy!20, rectangle, rounded corners=2pt, minimum height=0.85cm, minimum width=2.0cm, align=center] (sic) at (-1.6,2.0) {\textbf{SiC衬底}\\ 天岳先进、天科合达};
\node[draw, fill=navy!20, rectangle, rounded corners=2pt, minimum height=0.85cm, minimum width=2.0cm, align=center] (equipment) at (1.6,2.0) {\textbf{制造设备}\\ 中微公司、北方华创};

\node[font=\bfseries\color{navy}, fill=white, inner sep=2pt, anchor=south] at (-1.6, 2.7) {上游：材料与设备（弹性\ding{72}\ding{72}\ding{72}）};

% ===== 中游：芯片制造 =====
\node[draw, fill=accentblue!30, rectangle, rounded corners=2pt, minimum height=0.95cm, minimum width=2.8cm, align=center] (idm1) at (-4.8,-0.5) {\textbf{IDM龙头}\\ 华润微、士兰微\\ 时代电气、扬杰科技};
\node[draw, fill=accentblue!30, rectangle, rounded corners=2pt, minimum height=0.95cm, minimum width=2.8cm, align=center] (design) at (-1.6,-0.5) {\textbf{设计龙头}\\ 斯达半导、新洁能\\ 东微半导、芯朋微};
\node[draw, fill=accentblue!30, rectangle, rounded corners=2pt, minimum height=0.95cm, minimum width=2.8cm, align=center] (foundry) at (1.6,-0.5) {\textbf{代工制造}\\ 华虹半导体\\ 中芯国际、三安光电};

\node[font=\bfseries\color{accentblue}, fill=white, inner sep=2pt, anchor=south] at (-1.6, 0.22) {中游：芯片制造（弹性\ding{72}\ding{72}\ding{72}\ding{72}\ding{72}）};

% ===== 下游：应用 =====
\node[draw, fill=riskgreen!20, rectangle, rounded corners=2pt, minimum height=0.85cm, minimum width=2.0cm, align=center] (auto) at (-4.8,-3.0) {\textbf{新能源车}\\ 比亚迪、特斯拉};
\node[draw, fill=riskgreen!20, rectangle, rounded corners=2pt, minimum height=0.85cm, minimum width=2.0cm, align=center] (ai) at (-1.6,-3.0) {\textbf{AI/数据中心}\\ 英伟达、华为};
\node[draw, fill=riskgreen!20, rectangle, rounded corners=2pt, minimum height=0.85cm, minimum width=2.0cm, align=center] (pv) at (1.6,-3.0) {\textbf{光伏储能}\\ 阳光电源、华为};
\node[draw, fill=riskgreen!20, rectangle, rounded corners=2pt, minimum height=0.85cm, minimum width=2.0cm, align=center] (ind) at (4.8,-3.0) {\textbf{工业控制}\\ 汇川技术};

\node[font=\bfseries\color{riskgreen}, fill=white, inner sep=2pt, anchor=south] at (0, -2.3) {下游：应用（需求验证\ding{72}\ding{72}\ding{72}\ding{72}\ding{72}）};

% ===== 箭头 =====
\draw[->, thick, navy!60] (wafer.south) -- (idm1.north);
\draw[->, thick, navy!60] (sic.south) -- (design.north);
\draw[->, thick, navy!60] (equipment.south) -- (foundry.north);

\draw[->, thick, accentblue!60] (idm1.south) -- (auto.north);
\draw[->, thick, accentblue!60] (design.south) -- (ai.north);
\draw[->, thick, accentblue!60] (foundry.south) -- (pv.north);

\end{tikzpicture}%
}
\vspace{0.2cm}
\sourcenote{本研究团队整理}
\end{exhibitbox}

\subsection{关键标的}

基于业绩确定性、产业链卡位、估值安全边际和催化弹性四大维度，我们筛选出功率半导体板块五大核心标的：

\begin{exhibitbox}[表4-2：功率半导体核心标的一览]
\centering
\small
\setlength{\aboverulesep}{2pt}
\setlength{\belowrulesep}{0pt}
\begin{tabularx}{\linewidth}{L{2.8cm} c l >{\raggedright\arraybackslash}X}
\toprule
\textbf{公司名称} & \textbf{代码} & \textbf{核心定位} & \textbf{投资逻辑} \\
\midrule
\rowcolor{riskgreen!5}
时代电气 & 688187 & 估值洼地的IDM龙头 & 车规IGBT/SiC双轮驱动，估值仍按轨交装备定价，半导体业务占比提升带动估值重估 \\
\rowcolor{riskgreen!5}
斯达半导 & 603290 & IGBT国产替代龙头 & 车规IGBT份额快速提升，SiC进展领先，受益于新能源车和光伏储能需求增长 \\
\rowcolor{accentblue!5}
华润微 & 688396 & 全品类IDM平台 & 产品布局最全的本土功率半导体IDM，制造平台优势显著，MOS/IGBT/SiC全面发力 \\
\rowcolor{accentblue!5}
士兰微 & 600460 & 12英寸IDM先锋 & 12英寸产能持续释放，产品结构升级，IGBT和SiC业务进入收获期 \\
\rowcolor{accentblue!5}
扬杰科技 & 300373 & 盈利稳健的功率IDM & 二极管/MOS龙头，SiC布局领先，盈利能力强，估值相对合理 \\
\bottomrule
\end{tabularx}
\sourcenote{本研究团队整理}
\end{exhibitbox}

\textbf{时代电气（688187）}——估值最低的功率半导体IDM。公司是国内少数具备IGBT芯片设计、制造、封装测试全流程能力的IDM厂商，车规级IGBT和SiC进展领先。截至2026年6月，公司PE(TTM)约20倍，显著低于行业平均60--80倍的水平，估值折价主要源于市场将其视为轨交装备公司。随着功率半导体业务占比提升，估值修复空间巨大。

\vspace{0.2em}

\textbf{斯达半导（603290）}——IGBT国产替代龙头。公司是国内IGBT领域的领军企业，车规级IGBT已批量交付多家头部车企，第七代车规IGBT大规模装车，第八代芯片研发顺利推进。SiC MOSFET已推出样品并在客户端验证，有望在2026--2027年贡献收入增量。

\vspace{0.2em}

\textbf{华润微（688396）}——全品类IDM平台。公司是国内产品布局最全的功率半导体IDM厂商，覆盖MOSFET、IGBT、二极管、功率IC等全系列产品。制造平台优势显著，8英寸和12英寸产能稳步扩张。SiC和GaN第三代半导体布局领先，长期成长空间广阔。

\vspace{0.2em}

\textbf{士兰微（600460）}——12英寸IDM先锋。公司12英寸产线持续释放产能，产品结构持续升级，IGBT模块和SiC器件业务进入收获期。公司在车规级领域布局深厚，多款产品通过AEC-Q100/Q101认证。随着产能利用率提升和产品结构优化，盈利弹性较大。

\vspace{0.2em}

\textbf{扬杰科技（300373）}——盈利稳健的功率IDM。公司是国内二极管和低压MOS龙头，盈利能力行业领先。SiC二极管和MOSFET布局领先，已实现批量出货。公司财务质量优异，现金流充沛，估值相对合理，具备较高的安全边际。

\subsection{风险因素}

\begin{riskbox}[功率半导体板块主要风险因素]
\small
\setlength{\aboverulesep}{0pt}
\setlength{\belowrulesep}{0pt}
\begin{tabularx}{\linewidth}{X L{2.8cm} c L{4.5cm}}
\toprule
\textbf{风险类别} & \textbf{风险事件} & \textbf{影响程度} & \textbf{关注信号} \\
\midrule
需求端 & AI资本开支放缓 & \textcolor{riskred}{高} & 科技大厂CapEx指引下修、服务器出货不及预期 \\
需求端 & 新能源车销量不及预期 & \textcolor{riskred}{高} & 月度新能源车销量、车企排产计划下修 \\
供给端 & 行业产能过剩加剧 & \textcolor{riskamber}{中高} & 晶圆厂产能利用率、行业扩产计划调整 \\
竞争端 & 海外厂商价格战 & \textcolor{riskamber}{中高} & 英飞凌/安森美等龙头价格策略变化 \\
技术端 & SiC/GaN迭代不及预期 & \textcolor{riskamber}{中} & 良率提升进度慢于预期、量产时间推迟 \\
政策端 & 出口管制加码 & \textcolor{riskred}{高} & 美欧日半导体管制政策更新 \\
\bottomrule
\end{tabularx}
\end{riskbox}

\textbf{下行空间测算}：若行业需求复苏低于预期且价格竞争加剧，板块估值可能从当前50--60倍PE回落至35--40倍，叠加业绩下修，板块整体回调空间约25--35\%。但龙头公司凭借技术壁垒和客户优势，下行空间相对有限，约15--25\%。建议在回调中逢低布局优质龙头。

\section{电力公用事业：电价机制改革+AI用电增量，价值重估进行时}

电力公用事业长期以来被视为低增速、低波动、高股息的防御性板块，市场普遍给予10--15倍PE的"公用事业估值"。但我们认为，AI算力革命带来的用电增量爆发与电价机制市场化改革正在深刻改变电力行业的盈利模式与估值逻辑，板块正处于从"类债属性"向"成长属性"切换的关键拐点，价值重估空间巨大。

\subsection{景气度三维验证}

\textbf{业绩维度：用电量超预期增长，盈利质量持续改善}。2025年全国全社会用电量同比增长约8.5\%，显著高于市场年初预期的5--6\%。其中，第二产业用电量增长约6.5\%，第三产业用电量增长约12\%，城乡居民生活用电量增长约9\%。2026年Q1用电量同比增长9.2\%，延续快速增长态势。

盈利方面，2025年A股电力板块归母净利润同比增长约15\%，其中火电板块受益于煤价回落同比增长约30\%，水电板块来水偏丰业绩稳健，核电板块保持两位数增长。进入2026年，随着电价市场化改革推进和用电需求持续增长，电力企业盈利能力有望进一步提升。

\vspace{0.3em}

\textbf{产业维度：AI数据中心用电需求爆发，增量空间巨大}。AI算力的快速扩张正在带来电力需求的指数级增长。一个中型数据中心（10万个标准机柜）的年用电量约相当于一个中等城市的居民用电量。以GB200整机柜为例，单机柜功率可达70--80kW，是传统机柜的7--8倍。

我们测算，2025年全国数据中心用电量约为2,800--3,000亿千瓦时，占全社会用电量的约3.5\%。到2030年，数据中心用电量有望达到8,000--10,000亿千瓦时，占全社会用电量的比例提升至8--10\%。AI数据中心是电力需求增长最快的细分领域，年均增速超过20\%。

\begin{exhibitbox}[表4-3：数据中心电力需求增长测算]
\centering
\small
\setlength{\tabcolsep}{6pt}
\begin{tabular}{l r r r r}
\toprule
\textbf{指标} & \textbf{2023} & \textbf{2025E} & \textbf{2027E} & \textbf{2030E} \\
\midrule
数据中心机柜数（万架） & $\sim$650 & $\sim$900 & $\sim$1,250 & $\sim$1,800 \\
单机架平均功率（kW） & $\sim$4.5 & $\sim$5.5 & $\sim$7.0 & $\sim$9.0 \\
数据中心年用电量（亿kWh） & $\sim$2,200 & $\sim$2,900 & $\sim$4,500 & $\sim$7,500 \\
占全社会用电量比例 & $\sim$2.8\% & $\sim$3.5\% & $\sim$5.0\% & $\sim$7.8\% \\
\midrule
AI数据中心用电量（亿kWh） & $\sim$300 & $\sim$700 & $\sim$1,800 & $\sim$4,000 \\
AI数据中心用电占比 & $\sim$14\% & $\sim$24\% & $\sim$40\% & $\sim$53\% \\
\bottomrule
\end{tabular}
\sourcenote{工信部，中国电力企业联合会，本研究团队测算}
\end{exhibitbox}

\vspace{0.3em}

\textbf{政策维度：电价机制改革深化，盈利模式重塑}。2025年以来，电价市场化改革进程明显加快。一是燃煤发电上网电价市场化机制进一步完善，浮动区间扩大；二是容量电价机制在全国范围推广，火电从"电量型"向"容量+电量"双轨制转型；三是分时电价机制持续优化，峰谷价差扩大，尖峰电价机制逐步建立；四是数据中心、充电桩等新型负荷参与电力市场的政策不断完善。

此外，新型电力系统建设、虚拟电厂政策、源网荷储一体化等政策持续推进，为电力行业带来新的商业模式和盈利增长点。

\subsection{预期差来源分析}

当前市场对电力公用事业的认知存在三方面核心预期差：

\vspace{0.3em}

\textbf{预期差一：仍按传统公用事业估值，未充分定价AI用电增量}。市场普遍将电力股视为"类债"的防御性资产，给予10--15倍PE的估值。但AI数据中心带来的用电增量具有高增长、高持续性的特点，与传统经济周期下的用电增长有本质区别。若AI算力投资持续超预期，电力需求增速可能长期维持在6--8\%的水平，远超历史平均3--5\%的增速。我们认为，具备区位优势（靠近算力枢纽）和电源结构优势（绿电占比高）的电力企业，其AI相关业务应给予更高的成长估值。

\vspace{0.3em}

\textbf{预期差二：低估电价机制改革对盈利模式的重塑效应}。市场对电价改革的认识仍停留在"电价小幅上浮"的层面，未充分认识到容量电价、现货市场、辅助服务等新机制对电力企业盈利模式的深刻改变。以容量电价为例，全国容量电价机制全面推开后，火电企业的容量收入占比有望达到20--30\%，盈利稳定性大幅提升，估值中枢有望上移。同时，峰谷价差扩大和现货市场价格发现功能完善，也为调峰能力强的电源带来额外收益。

\vspace{0.3em}

\textbf{预期差三：新能源消纳压力被过度担忧}。市场普遍担心新能源大规模接入后的消纳问题，认为弃风弃光率将持续攀升。但实际上，特高压输电通道建设、储能装机快速增长、电力市场机制完善、需求侧响应能力提升等多方面因素正在系统性改善新能源消纳条件。2025年全国风电、光伏发电利用率分别达到97\%和98\%，处于较高水平。我们认为，新能源消纳压力将呈现结构性、区域性特征，而非全面性紧张，优质新能源运营商的盈利能力有保障。

\subsection{核心催化与时间窗口}

未来12--18个月，电力公用事业板块存在四大核心催化：

\vspace{0.3em}

\textbf{催化一：容量电价机制全国推广（2026年下半年--2027年）}。容量电价机制从试点省份向全国推广，火电企业盈利模式重塑，从单纯赚取电量收益转向"容量收益+电量收益+辅助服务收益"的多元模式。关注指标：各省容量电价政策落地情况、容量电价标准、火电盈利变化。

\vspace{0.3em}

\textbf{催化二：虚拟电厂政策密集落地（2026年下半年）}。虚拟电厂、需求侧响应、电力辅助服务等政策持续推进，为电网企业和综合能源服务商带来新的增长点。关注指标：虚拟电厂试点规模、辅助服务市场交易电量、相关政策文件出台。

\vspace{0.3em}

\textbf{催化三：数据中心用电需求持续超预期（2026H2--2027）}。AI算力投资持续高增长，带动数据中心用电需求超预期，电力供需格局趋紧，电价易涨难跌。关注指标：数据中心新增装机、AI服务器出货量、各省份用电增速数据。

\vspace{0.3em}

\textbf{催化四：电力现货市场全面运行（2026--2027）}。电力现货市场在全国范围内全面运行，价格发现功能完善，峰谷价差进一步扩大，调峰价值凸显。关注指标：现货市场运行省份数量、峰谷价差水平、现货电价走势。

\subsection{产业链图谱与受益环节}

电力产业链可分为发电、输配电、售电及用电服务三个环节。AI用电增量和电价改革背景下，各环节均有不同程度的受益。

\begin{exhibitbox}[图4-2：电力产业链与AI用电增量受益环节]
\centering
\vspace{0.6cm}
\makebox[\textwidth][c]{%
\begin{tikzpicture}[node distance=0.5cm, every node/.style={font=\small}]

% 发电侧
\node[draw, fill=navy!25, rectangle, rounded corners=2pt, minimum height=0.85cm, minimum width=2.0cm, align=center] (nuclear) at (-5.25, 3.0) {\textbf{核电}\\ 中国广核、中国核电};
\node[draw, fill=navy!25, rectangle, rounded corners=2pt, minimum height=0.85cm, minimum width=2.0cm, align=center] (hydro) at (-1.75, 3.0) {\textbf{水电}\\ 长江电力、华能水电};
\node[draw, fill=navy!25, rectangle, rounded corners=2pt, minimum height=0.85cm, minimum width=2.0cm, align=center] (thermal) at (1.75, 3.0) {\textbf{火电}\\ 华能国际、宝新能源};
\node[draw, fill=navy!25, rectangle, rounded corners=2pt, minimum height=0.85cm, minimum width=2.0cm, align=center] (renewable) at (5.25, 3.0) {\textbf{新能源}\\ 三峡能源、龙源电力};

\path (nuclear.west) -- (renewable.east) node[midway, above=0.8cm, font=\bfseries\color{navy}] {发电侧（受益\ding{72}\ding{72}\ding{72}\ding{72}）};

% 输配电
\node[draw, fill=accentblue!30, rectangle, rounded corners=2pt, minimum height=0.95cm, minimum width=6.5cm, align=center] (grid) at (0, 0.8) {\textbf{输配电与电网自动化}\\ 国家电网/南方电网体系 · 国电南瑞、许继电气、南瑞继保};

\node[above=0.5cm of grid, font=\bfseries\color{accentblue}] {输配电侧（受益\ding{72}\ding{72}\ding{72}）};

% 售电/服务
\node[draw, fill=riskgreen!20, rectangle, rounded corners=2pt, minimum height=0.85cm, minimum width=2.5cm, align=center] (vpp) at (-1.9, -1.4) {\textbf{虚拟电厂/综合能源}\\ 国电南瑞、国网信通};
\node[draw, fill=riskgreen!20, rectangle, rounded corners=2pt, minimum height=0.85cm, minimum width=2.5cm, align=center] (retail) at (1.9, -1.4) {\textbf{售电公司}\\ 地方能源集团};

\path (vpp.west) -- (retail.east) node[midway, above=0.8cm, font=\bfseries\color{riskgreen}] {售电与用电服务（受益\ding{72}\ding{72}\ding{72}\ding{72}）};

% 下游
\node[draw, fill=riskamber!25, rectangle, rounded corners=2pt, minimum height=0.85cm, minimum width=3.0cm, align=center] (dc) at (0, -3.0) {\textbf{AI数据中心}\\ 运营商 · 云厂商 · 算力企业};

\node[below=0.4cm of dc, font=\bfseries\color{riskamber}] {需求侧：用电增量引擎（驱动\ding{72}\ding{72}\ding{72}\ding{72}\ding{72}）};

% 箭头（发电→电网，均匀分布在grid顶部）
\draw[->, thick, navy!60] (nuclear.south) -- ([xshift=-2.0cm]grid.north);
\draw[->, thick, navy!60] (hydro.south) -- ([xshift=-0.7cm]grid.north);
\draw[->, thick, navy!60] (thermal.south) -- ([xshift=0.7cm]grid.north);
\draw[->, thick, navy!60] (renewable.south) -- ([xshift=2.0cm]grid.north);

\draw[->, thick, accentblue!60] (grid.south) -- (vpp.north);
\draw[->, thick, accentblue!60] (grid.south) -- (retail.north);

\draw[->, thick, riskgreen!60] (vpp.south) -- (dc.north);
\draw[->, thick, riskgreen!60] (retail.south) -- (dc.north);

\end{tikzpicture}%
}
\vspace{0.2cm}
\sourcenote{本研究团队整理}
\end{exhibitbox}

\subsection{关键标的}

基于资源禀赋、区位优势、盈利弹性和估值水平四大维度，我们筛选出电力公用事业板块五大核心标的：

\begin{exhibitbox}[表4-4：电力公用事业核心标的一览]
\centering
\small
\setlength{\aboverulesep}{2pt}
\setlength{\belowrulesep}{0pt}
\begin{tabularx}{\linewidth}{L{2.8cm} c l >{\raggedright\arraybackslash}X}
\toprule
\textbf{公司名称} & \textbf{代码} & \textbf{核心定位} & \textbf{投资逻辑} \\
\midrule
\rowcolor{riskgreen!5}
长江电力 & 600900 & 水电龙头，价值压舱石 & 全球最大水电上市公司，现金流稳定，分红率高，兼具防御性与成长属性 \\
\rowcolor{riskgreen!5}
国电南瑞 & 600406 & 电网自动化绝对龙头 & 新型电力系统建设核心受益者，虚拟电厂和AI调度技术领先，估值有望提升 \\
\rowcolor{accentblue!5}
三峡能源 & 600905 & 新能源运营龙头 & 风光资源优质，装机快速增长，绿电交易溢价，受益于AI算力绿电需求 \\
\rowcolor{accentblue!5}
华能国际 & 600011 & 火电龙头，弹性最大 & 火电装机规模领先，受益于容量电价和用电需求增长，盈利弹性大 \\
\rowcolor{accentblue!5}
宝新能源 & 000690 & 区域火电龙头 & 广东区域火电龙头，区位优势显著，电价市场化弹性大，估值偏低 \\
\bottomrule
\end{tabularx}
\sourcenote{本研究团队整理}
\end{exhibitbox}

\textbf{长江电力（600900）}——水电龙头，价值压舱石。公司是全球最大的水电上市公司，拥有三峡、葛洲坝、溪洛渡、向家坝、乌东德、白鹤滩六座巨型水电站，总装机容量约7,170万千瓦。公司现金流充沛，分红率长期维持在70\%左右，股息率约3.5--4\%，是稳健配置的首选。同时，来水改善和电价上浮也为业绩提供弹性。

\vspace{0.2em}

\textbf{国电南瑞（600406）}——电网自动化绝对龙头。公司是国家电网旗下核心上市平台，在电网调度、继电保护、变电站自动化等领域市占率领先。新型电力系统建设、特高压投资、虚拟电厂发展均为公司带来持续增长机遇。公司在AI+电网调度领域技术储备深厚，有望受益于AI与电力融合的大趋势。

\vspace{0.2em}

\textbf{三峡能源（600905）}——新能源运营龙头。公司是三峡集团旗下新能源运营平台，风电和光伏装机规模快速增长，资源禀赋优质。公司积极布局源网荷储一体化和制氢等新业务，长期成长空间广阔。AI数据中心对绿电需求的增长，也为公司带来新的发展机遇。

\vspace{0.2em}

\textbf{华能国际（600011）}——火电龙头，弹性最大。公司是国内装机规模最大的火电上市公司之一，火电装机超1.2亿千瓦。公司将显著受益于容量电价机制推广和用电需求增长，盈利弹性较大。同时，公司新能源装机快速增长，电源结构持续优化。

\vspace{0.2em}

\textbf{宝新能源（000690）}——区域火电龙头。公司是广东省重要的电力供应商，区位优势显著。广东是全国电力供需最紧张、电价市场化程度最高的省份之一，公司将受益于电价上浮和峰谷价差扩大。公司估值偏低，股息率较高，具备较好的安全边际。

\subsection{风险因素}

\needspace{6\baselineskip}
\begin{riskbox}[电力公用事业板块主要风险因素]
\small
\setlength{\aboverulesep}{0pt}
\setlength{\belowrulesep}{0pt}
\begin{tabularx}{\linewidth}{X L{2.8cm} c L{4.5cm}}
\toprule
\textbf{风险类别} & \textbf{风险事件} & \textbf{影响程度} & \textbf{关注信号} \\
\midrule
成本端 & 煤炭价格大幅上涨 & \textcolor{riskred}{高} & 环渤海动力煤价格指数、进口煤价 \\
政策端 & 电价改革不及预期 & \textcolor{riskred}{高} & 电价政策文件、浮动区间调整力度 \\
需求端 & 宏观经济下行用电疲软 & \textcolor{riskamber}{中高} & 月度用电量数据、工业增加值 \\
消纳端 & 新能源消纳压力加大 & \textcolor{riskamber}{中} & 弃风弃光率、利用小时数变化 \\
气候端 & 极端天气来水异常 & \textcolor{riskamber}{中} & 降水量、水库水位、气候预测 \\
\bottomrule
\end{tabularx}
\end{riskbox}

\textbf{下行空间测算}：若煤炭价格大幅反弹或电价改革推进缓慢，火电企业盈利可能承压，板块估值可能从当前12--15倍PE回落至8--10倍，叠加业绩下修，火电板块回调空间约25--35\%。水电和新能源运营商由于现金流稳定、分红率高，下行空间相对有限，约10--20\%。建议关注煤价走势和政策落地节奏，逢低布局优质标的。

\vspace{0.5em}

\needspace{8\baselineskip}
\begin{keyinsight}[本章核心观点]
\textbf{功率半导体}：AI算力革命与能源转型双轮驱动，国产替代从消费级向车规级纵深推进。市场低估AI侧增量弹性和SiC/GaN渗透速度，板块处于复苏初期，配置价值突出。重点关注时代电气、斯达半导、华润微、士兰微、扬杰科技。

\vspace{0.3em}

\textbf{电力公用事业}：AI用电增量爆发与电价机制改革共振，板块正从"类债"防御属性向"成长+价值"属性切换，价值重估空间巨大。市场仍按传统公用事业估值，未充分定价AI用电增量和盈利模式重塑。重点关注长江电力、国电南瑞、三峡能源、华能国际、宝新能源。
\end{keyinsight}
